% ============ 1. Introduction ============
\section{Introduction}\label{sec:intro}

Large language models have moved from curiosity to instrument in empirical finance. \citet{lopezlira2026}, in the \emph{Journal of Financial Economics}, score firm-specific news headlines with ChatGPT and show that the scores predict next-day stock returns and that this predictability is not matched by dictionary-based tone measures. The finding has been replicated, extended, and contested across models, markets, and windows \citep{nie2024survey}, and it sits at the center of a rapidly growing literature that treats LLM output as a measurement layer for text-based variables. The premise of that literature is that the LLM \emph{reads the text}. The premise is testable, and the tests matter: if the scores partly reflect the model's knowledge about the \emph{firm} rather than the \emph{headline}, or the \emph{already-realized price move} described in retrospective headlines, or the \emph{luck of a single stochastic call}, or the \emph{selection of which firm-days carry news at all}, then predictability statistics confound semantics with contamination. \citet{glasserman2023} show for in-sample windows that company identity itself distorts GPT sentiment scores enough to change backtest performance; \citet{lopezlira2025} show that LLMs memorize economic and financial data inside their training window; \citet{didisheim2025} show the memory is stronger for indices than for individual stocks; \citet{chen2025} show that LLMs extrapolate recent returns when asked to forecast; and \citet{herrera2025} document large cross-model and repeat-call variability in LLM sentiment classifications. What no study has done is measure these channels \emph{jointly}, on the \emph{same} headlines, against the \emph{same} return windows, under a common inference standard. That is what this paper does.

Our vehicle is deliberately modest: a firm-level panel of 36 large U.S.\ equities over 74 trading days (July 1 to October 14, 2026), 542 firm-specific headlines collected through a documented web-search protocol, of which 341 are dated inside the window and 240 relevance-qualified firm-days survive aggregation. The window postdates the plausible training cutoff of the scoring model, which removes the most literal form of look-ahead bias~---memorized realized outcomes \citep{lopezlira2025,didisheim2025}~---and therefore isolates the subtler channels that remain operative out of sample. The same 341 headlines are scored four ways: (i)~the original protocol, with firm identity visible; (ii)~an \emph{anonymized} protocol in which company names and tickers are replaced by ``the company''; (iii)~three \emph{independent repeat calls} of the original protocol; and (iv)~the \citet{loughran2011} dictionary. Each condition changes exactly one ingredient of the measurement. We then pre-specify a family of five hypotheses~---baseline predictability, identity effects, scoring stability, price-echo contamination, and coverage selection~---and evaluate all of them under two-way clustered inference, wild-cluster bootstrap, day-cluster bootstrap difference tests, and Benjamini--Hochberg control of the false discovery rate over the declared family.

Five results summarize the audit.

\emph{First, the baseline result is real but fragile~---and its significance is draw-dependent.} Next-day returns regress on named LLM sentiment with $t=2.10$ ($p=0.035$) under two-way clustered errors, and the wild-cluster bootstrap confirms the signal is not an artifact of a handful of clusters ($p=0.030$). But the same statistic ranks last in the family once the Benjamini--Hochberg correction is applied ($q=0.062$), the coefficient loses conventional significance when the same-day return~---which is itself strongly correlated with sentiment ($t=3.3$)~---is added as a control ($t=1.82$), and re-scoring the identical headlines in thirteen independent passes across five sittings moves the $t$-statistic between $1.47$ and $2.12$: the point estimate is stable, but whether a single-call pipeline would have reported ``significant'' depends on the day it scored. In a sample of 240 firm-days, the honest summary is ``marginal evidence,'' not ``significant predictability.''

\emph{Second, what looked like an identity prior was a session artifact~---and unpicking the artifact exposed two real level channels.} Removing firm identifiers reduces the average score by 4.75 points (paired $t=5.14$)~---a level shift that first read as the out-of-sample analog of the ``distraction effect'' of \citet{glasserman2023}. A same-session benchmark overturned that reading: the named and anonymized rounds were fielded twelve hours apart, and the identical named protocol drifted $+5.79$ points between sittings ($t=7.03$), dwarfing a same-session residual of roughly one point. The experiment that the diagnosis demanded~---a third sitting with four named and four anonymized rounds interleaved at the batch level, pre-specified before data collection~---then settled the question in both directions at once. The interleaved identity effect is $-2.82$ points (all four round-level estimates negative; round-level $t=6.8$; batch-call-clustered $t=2.7$; V-excluded $-3.13$ at $t=4.8$): anonymized scores sit systematically \emph{above} their named twins, a real but mild de-identification effect, opposite in sign to the two-session shift. And the pre-specified swap control of a fifth session settles what the effect is made of: replacing each headline's firm identifiers with a \emph{different} firm's leaves levels essentially unchanged ($+0.34$, $p=0.69$), while the placeholder rewrite itself moves them by two-plus points~---the level effect is the removal of a concrete name, not a prior about the identified firm. And the drift itself decomposed: re-scoring the identical headlines inside the original 542-article batches (same session, same endpoint) moves levels $+3.61$ points, replicated at $+3.50$ in a fourth sitting with the run order inverted (pooled $z=3.4$)~---the ``session drift'' is mostly a \emph{batch-composition} effect, a measurement channel we believe has not previously been quantified in this setting, with a further twist: random re-batching of the same corpus moves levels too, while the linear part of the anchoring we estimate is provably level-neutral for re-batching and quantifiably explains only about 1\% of the composition shift~---the batch channel lives in the response's non-linear components. What identity does \emph{not} do is equally clear and now carries a formal equivalence verdict: the return-predictive coefficient is \emph{formally equivalent} between the named and anonymized conditions at a margin of half the benchmark coefficient in every comparison~---cross-session, within-session, interleaved, and the corrected fifth-session variant, with 90\% intervals closing at 21--42\% of the benchmark~---while the swap variant is the single near-miss (its interval needs a 56\% margin) and the dictionary horse race settles nothing in either direction (its interval spans $4.4\times$ the benchmark). The lesson is uncomfortable on both sides of the LLM-versus-dictionary debate: in a sample this size, the difference between instruments is not estimable, and any claim of dominance based on one instrument being significant while the other is not is a statistical error we document explicitly below.

\emph{Third, the instrument is stochastic.} Four independent calls of the \emph{same} prompt on the \emph{same} headlines produce an intra-class correlation of $0.919$ (95\% CI $[0.90, 0.93]$) and only a 29.9\% exact-replication rate across all four calls (44\% for a single pair of calls); the within-headline standard deviation across calls is 7.8 points. This contradicts the intuition that default decoding makes the instrument deterministic~---an intuition we initially shared as well~---and temperature-zero decoding does not restore determinism either: two identical temperature-zero rounds agree bitwise on only 83.3\% of headlines, so part of the noise is server-side. Every downstream regression inherits this noise as attenuation.

\emph{Fourth, the predictive content directionally sits where the text echoes prices.} A text audit flags 15.2\% of headlines as containing language about realized price moves (``rose 14\%,'' ``dropped by 5\% after earnings''). The next-day coefficient on sentiment is $0.047$ ($t=1.80$) for firm-days with such headlines and $0.003$ ($t=0.61$) for firm-days without them; the formal difference is directional ($p=0.10$) rather than established~---underpowered in the tails~---so we treat the echo channel as a well-corroborated direction, not a certified finding. The anonymized condition does essentially nothing to the same-day loading of sentiment on returns ($t=3.33$ named vs.\ $3.11$ anonymized), which rules out the scorer's memory as the source: the echo is \emph{in the text}. Because our window postdates the training cutoff, this channel is precisely the one the literature cannot dismiss as memorization~---search-collected news archives systematically over-represent retrospective, price-descriptive coverage.

\emph{Fifth, the sample itself is selected.} Covered firm-days exhibit mean absolute same-day returns of 215 \bps{} against 154 \bps{} for uncovered firm-days (firm-and-day fixed-effects test $t=2.76$; see Section~\ref{sec:results:h5} for why the pooled two-sample statistic overstates this), and firm-level news counts correlate 0.48 with firm return volatility. The estimation sample is therefore an over-weight of high-attention, high-absolute-return events~---exactly the environment in which short-horizon continuation and reversal are strongest for reasons unrelated to news semantics \citep{chan2003}.

None of these channels is large enough, alone, to overturn a well-powered result. Jointly, they change what a small-sample LLM predictability statistic \emph{means}. The contribution of this paper is therefore not a new anomaly but a \emph{measurement audit protocol}: anonymization, repeat scoring, text-level echo flags, and coverage-selection diagnostics, all run on one frozen corpus under one inference standard~---together with the staged, pre-specified re-fielding sessions that repair the confounds the two-session fielding created. We release the full pipeline~---headlines, all archived score sets, prompts, scripts, and every table and figure~---as a replication package, and we report the family-wise error control explicitly so that future studies can register audit families before examining returns.

The remainder proceeds as follows. Section~\ref{sec:lit} positions the audit against the LLM-finance and textual-analysis literature. Section~\ref{sec:framework} sets out the measurement decomposition that organizes the tests. Section~\ref{sec:data} documents the corpus, the four scoring conditions, and the econometric design. Section~\ref{sec:results} presents the five audit results. Section~\ref{sec:robust} reports robustness and the threats to validity we could and could not address. Section~\ref{sec:discussion} discusses what survives, for whom, and under which protocol. Section~\ref{sec:limitations} states limitations; Section~\ref{sec:conclusion} concludes.
